% =========================
% report/discussion.tex
% =========================

\subsection{Assumptions and Validity Domain}
The method assumes that the relevant template surface is planar (or sufficiently close for the required tolerance). When this assumption is violated (e.g., curved surfaces, significant depth variation), a single homography cannot model the mapping accurately, and residual errors will persist even with correct matching.

\subsection{Limitations}
Key limitations include:
\begin{itemize}
  \item Sensitivity to glare, low texture, and repetitive patterns.
  \item Degradation under large perspective changes or partial occlusions.
  \item Dependence on template quality (line thickness, contrast, and scale standardization).
\end{itemize}

\subsection{Mitigation Strategies (Evidence-Driven)}
Mitigations should be adopted only if they reduce tail failures and improve success rate under the chosen tolerance:
\begin{itemize}
  \item ROI gating to suppress background correspondences.
  \item Parameter selection guided by ablation stability regions.
  \item Template enhancement (contrast normalization; multi-variant templates when scale varies).
\end{itemize}

\subsection{Threats to Validity}
Threats include annotation noise (for reprojection metrics), dataset bias (limited variety of lighting/viewpoints), and over-tuning parameters on a small dataset. These are mitigated by reporting distributions, including failure cases, and separating tuning from final evaluation.

